\documentclass[11pt]{article}

\usepackage[a4paper,margin=29mm]{geometry}
\usepackage[T1]{fontenc}
\usepackage{lmodern}
\usepackage{microtype}
\usepackage{amsmath,amssymb,amsthm,mathtools}
\usepackage{booktabs}
\usepackage{array}
\usepackage{float}
\usepackage{enumitem}
\usepackage[hidelinks]{hyperref}

\newtheorem{theorem}{Theorem}[section]
\newtheorem{proposition}[theorem]{Proposition}
\newtheorem{corollary}[theorem]{Corollary}
\newtheorem{lemma}[theorem]{Lemma}
\theoremstyle{remark}
\newtheorem{remark}[theorem]{Remark}

\newcommand{\PG}{\mathrm{PG}}
\newcommand{\F}{\mathbb{F}}
\newcommand{\cC}{\mathcal{C}}
\newcommand{\cU}{\mathcal{U}}

\title{Computational strengthenings of Clebsch syndrome rigidity}
\author{Tavis Rudd}
\date{July 2026}

\begin{document}
\maketitle

\begin{abstract}
For a projective arc $A\subset\PG(2,q)$, let $\cU(A)$ be the points on no
chord of $A$.  The geometric paper proves, without an exhaustive
classification of six-arcs over $\F_{11}$, that a six-arc in $\PG(2,11)$
whose uncovered locus lies on a conic is the Clebsch hexagon.  Here exact
finite computation sharpens and extends that result.
There are fifteen projective classes of six-arcs over $\F_{11}$; the
Clebsch class is the unique one whose uncovered locus is contained in a
cubic, and it is separated from every other class by a four-point gap in
uncovered-set size.  A Sylvester-graph obstruction shows that $q=11$ is the
only field order admitting a conic-filling six-arc.  Exhaustive orbit
searches then classify all conic-filling arcs through eight points: only
the projective four-frame over $\F_5$ and the Clebsch six-arc over
$\F_{11}$ occur.
At the terminal field $q=13$, Segre tangent triples give a conceptual
exclusion of the saturated eight-point case.  The associated binary
passant/internal incidence code has exact minimum distance $12$; its
$364$ minimum words reconstruct the passant incidence matrix and its full
$\operatorname{PGL}(2,13)$ symmetry.
The finite claims are accompanied by exact replay routes and a
claim-by-claim trust ledger.
\end{abstract}

\noindent\textbf{Keywords.} finite projective plane; arc; conic; MDS code;
deep hole; exhaustive classification.

\smallskip
\noindent\textbf{2020 Mathematics Subject Classification.} 51E21, 94B05,
51E22, 05B25.

\section{Setup and relation to the geometric paper}

A \emph{$k$-arc} is a set $A$ of $k$ points of $\PG(2,q)$ with no three
collinear.  Its chords are the joins of pairs of vertices, and
\[
 \cU(A)=\{P\in\PG(2,q): P\text{ lies on no chord of }A\}.
\]
We call $A$ \emph{conic-filling} when $\cU(A)$ is the full rational point
set of a nonsingular conic.  If $n_i$ counts the off-arc points incident
with exactly $i$ chords, put
\[
 r=\lfloor k/2\rfloor,\qquad
 t(A)=\sum_{i=3}^{r}\binom{i-1}{2}n_i,\qquad
 m=\binom{k}{2}.
\]

The geometric paper~\cite{RuddRigidity2026} proves the following three
inputs.  They are restated to make every dependency in the finite arguments
explicit.

\begin{theorem}[Chord defect {\cite[Theorem 3.1]{RuddRigidity2026}}]
\label{lem:chord-defect}
For every $k$-arc with $k\ge4$,
\[
 |\cU(A)|=q^2-(m-1)q+
 \left(m+3\binom{k}{4}-k+1\right)-t(A),
\qquad
0\le t(A)\le3\binom{k}{4}\frac{r-2}{r}.
\]
\end{theorem}

\begin{theorem}[Conic-filling window
{\cite[Theorem 3.2]{RuddRigidity2026}}]
\label{cor:conic-filling-window}
If $\cU(A)$ is a nonsingular conic, then $q$ is odd and
\[
 2k-3\le q\le\frac{k(k-1)+3}{3}<\binom{k}{2}.
\]
For $k=6$, writing $c(A)=n_3$ gives
\[
 |\cU(A)|=q^2-14q+55-c(A),\qquad
 0\le c(A)\le15,\qquad
 c(A)=(q-6)(q-9)
\]
when $|\cU(A)|=q+1$.
\end{theorem}

\begin{theorem}[Geometric rigidity
{\cite[Theorem 1.1]{RuddRigidity2026}}]
\label{thm:rigidity}
For a six-arc $A\subset\PG(2,11)$, the locus $\cU(A)$ lies on a conic if
and only if $A$ is projectively equivalent to the Clebsch hexagon.  In that
case $\cU(A)$ is exactly a nonsingular conic.
\end{theorem}

The columns of a parity-check matrix of a projective
$[k,k-3,4]_q$ MDS code form a $k$-arc.  When $\cU(A)$ is nonempty it is
the projective distance-three syndrome locus.  Projective equivalence of
the unordered column sets is the same as monomial equivalence of the
codes.  Thus each geometric classification below has the stated coding
interpretation.

\section{The six-arc census over \texorpdfstring{$\F_{11}$}{F11}}

\paragraph{Computer-assisted conic-inscribed subcensus.}
For comparison, the exact replay enumerates all $1548$ frame-normalized
six-arc presentations and selects the $252$ whose vertices lie on a
nonsingular conic. Their $|\cU(A)|$ histogram is
\[
\begin{array}{c|rrrr}
|\cU(A)|&18&19&20&22\\ \hline
\text{frequency}&30&60&90&72
\end{array}
\]
so
$|\cU(A)|\in\{18,19,20,22\}$. The replay is listed in the ``Fifteen
classes and numerical gap'' row of Table~\ref{tab:replays}.

Table~\ref{tab:fifteen-classes} records the complete census needed below.
Every representative
contains the standard frame
\[
 (1,0,0),\ (0,1,0),\ (0,0,1),\ (1,1,1);
\]
the table gives the remaining two points of each representative.
Here $d_{\min}$ is the least degree of a nonzero form vanishing on the
uncovered locus.

\begin{table}[H]
\centering
\scriptsize
\begin{tabular}{@{}c c r r r@{}}
\toprule
class & pair & $|G|$ & $|\cU|$ & $d_{\min}$\\
\midrule
\texttt{C01}&23/32& 2&20&5\\
\texttt{C02}&23/34& 6&18&4\\
\texttt{C03}&23/36& 3&19&5\\
\texttt{C04}&23/37& 3&19&5\\
\texttt{C05}&23/39& 4&20&5\\
\bottomrule
\end{tabular}
\hfill
\begin{tabular}{@{}c c r r r@{}}
\toprule
class & pair & $|G|$ & $|\cU|$ & $d_{\min}$\\
\midrule
\texttt{C06}&23/42& 2&20&5\\
\texttt{C07}&23/48& 1&21&5\\
\texttt{C08}&23/49& 4&20&5\\
\texttt{C09}&23/56& 4&20&5\\
\texttt{C10}&23/62& 6&19&5\\
\bottomrule
\end{tabular}
\hfill
\begin{tabular}{@{}c c r r r@{}}
\toprule
class & pair & $|G|$ & $|\cU|$ & $d_{\min}$\\
\midrule
\texttt{C11}&23/67&12&16&4\\
\texttt{C12}&23/74& 5&22&6\\
\texttt{C13}&23/89& 6&18&4\\
\texttt{C14}&23/(9,10)&12&18&4\\
\texttt{C15}&34/45&60&12&2\\
\bottomrule
\end{tabular}
\caption{The fifteen projective classes of six-arcs in $\PG(2,11)$.
In the pair column, \(xy/uv\) abbreviates
\((1,x,y),(1,u,v)\).
Class \texttt{C15} is the Clebsch class.  Both replay programs regenerate
the same canonical keys; the global-gap replay supplies $|G|$ and
$|\cU|$, and the low-degree replay supplies $d_{\min}$.}
\label{tab:fifteen-classes}
\end{table}

The same census supports a strengthening in which a cubic, rather than a
conic, is allowed.

\begin{proposition}[Low-degree algebraic rigidity]
\label{prop:low-degree-rigidity}
For a six-arc $A\subset\PG(2,11)$, the uncovered locus $\cU(A)$ is
contained in the zero set of a nonzero homogeneous form of degree at most
three if and only if $A$ is projectively equivalent to the Clebsch
hexagon. In the Clebsch case the least such degree is two: the quadratic
vanishing space is spanned by the defining nonsingular conic equation $Q$,
and the cubic vanishing space is
\[
 Q\cdot\F_{11}[X,Y,Z]_1.
\]
\end{proposition}

\begin{proof}[Computer-assisted proof]
For each of the fifteen projective classes in the six-arc census, evaluate
the ten cubic monomials on $\cU(A)$. The resulting evaluation map has rank
ten for every non-Clebsch class and rank seven for the Clebsch class. Thus
no nonzero cubic vanishes on a non-Clebsch uncovered locus. A containing
form of smaller positive degree could be multiplied by a nonzero monomial
to produce such a cubic, so smaller positive degrees are excluded as well;
no nonzero constant vanishes on the nonempty locus. For the
Clebsch class, the quadratic kernel is generated by $Q$, and multiplication
by linear forms gives the full three-dimensional cubic kernel.
\end{proof}

\begin{corollary}[Monomial characterization]
\label{cor:monomial-characterization}
Up to monomial equivalence, there is a unique linear $[6,3,4]_{11}$ code
of covering radius three whose projective deep-hole syndrome locus is
contained in a plane curve of degree at most three. It is the Clebsch code,
whose projective stabilizer is $A_5$.
\end{corollary}

\begin{proof}
The columns of a parity-check matrix of such a code form a six-arc $A$, and
covering radius three identifies its projective deep-hole syndrome locus
with $\cU(A)$. Proposition~\ref{prop:low-degree-rigidity} makes $A$
projectively equivalent to the Clebsch hexagon. A projectivity between two
unordered parity-check column sets lifts, after choosing column
representatives, to a row operation together with a coordinate permutation
and nonzero coordinate scalings; these are exactly a monomial code
equivalence. Conversely, the Clebsch code has the stated property, and
Dye's Theorem~3 gives the $A_5$ statement
\cite[p.~278]{Dye1991}.
\end{proof}

\begin{remark}[Sharpness of the degree threshold]
\label{rem:degree-threshold}
The cutoff cannot be raised from three to four: the four non-Clebsch
classes \texttt{C02}, \texttt{C11}, \texttt{C13}, and \texttt{C14} have
$d_{\min}=4$, so the failure is not isolated. For \texttt{C02}, the
low-degree replay in Table~\ref{tab:replays} prints a quartic whose rational
zero set is exactly the $18$-point uncovered locus and verifies the equality
point by point. No smoothness, irreducibility, or genus assertion is used
here.
\end{remark}

\begin{theorem}[Numerical gap]
\label{thm:gap}
The Clebsch class has $|\cU(A)|=12$, while every non-Clebsch
six-arc has $|\cU(A)|\ge16$. Equivalently, $|\cU(A)|\le15$
characterizes the Clebsch class, for which $|\cU(A)|=12$.
If $|\cU(A)|\ge16$, then $\cU(A)$ lies on no conic.
\end{theorem}

\begin{proof}[Computer-assisted proof]
Any ordered four points of a six-arc form a projective frame, uniquely
normalized to
$(1,0,0),(0,1,0),(0,0,1),(1,1,1)$. Enumerating the remaining unordered pair
subject to the arc condition gives $1548$ normalized sets. Each embedded arc
has $\binom{6}{4}\,4!=360$ ordered-frame choices. For each choice, apply the
unique normalizing projectivity, sort the resulting six normalized points,
and encode that list lexicographically. Taking the least of the $360$ lists
removes the frame choice; two arcs have the same least list exactly when a
projectivity carries one to the other. This canonical key gives fifteen
classes, and the number of distinct normalized presentations of a class
with stabilizer $G$ is $360/|G|$. Their extension counts are
\[
 |\cU(A)|\in\{12,16,18,19,20,21,22\},
\]
with multiplicities
\[
 (6,30,150,300,630,360,72)
\]
across the $1548$ normalized representatives. The six presentations
attaining $12$ form the single Clebsch class; every other class has at
least $16$ extension points. Theorem~\ref{thm:rigidity} supplies the final
conic-containment assertion.
\end{proof}

\section{Cross-field uniqueness for six-arcs}

\begin{lemma}[$q=9$ polarity graph]
\label{lem:q9-polarity}
For a nonsingular conic $\cC\subset\PG(2,9)$, let $\Sigma$ be the graph on
the $36$ internal points in which
\[
 P\sim Q\quad\Longleftrightarrow\quad Q\in P^\perp
\]
for the conic polarity. Then $\Sigma$ is the Sylvester graph, with
intersection array
\[
 \{5,4,2;1,1,4\}.
\]
For distinct internal points $P,Q$, the line $PQ$ is passant to $\cC$ if
and only if $P$ and $Q$ are at distance two in $\Sigma$, and the graph
joining pairs at distance two has clique number
\[
 \omega(\Sigma_2)=5.
\]
\end{lemma}

\begin{proof}
Counting internal points on polar lines and their intersections gives
$b_0=5$, $b_1=4$, $b_2=2$, $c_1=c_2=1$, and $c_3=4$, hence the displayed
intersection array. This is the Sylvester graph
\cite[\S13.1.2]{BrouwerCohenNeumaier1989}; uniqueness for this intersection
array is proved in~\cite[Thm.~6]{JurisicVidali2019}.

For distinct internal points $P,Q$, put
$R=P^\perp\cap Q^\perp=(PQ)^\perp$. If $PQ$ is passant, then its pole $R$
is internal and adjacent to both $P$ and $Q$; since the intersection array
has $a_1=0$, these vertices are at distance two. Conversely, a common
neighbor of $P$ and $Q$ must equal $R$, so distance two makes the pole of
$PQ$ internal and hence makes $PQ$ passant. The final value is
$\operatorname{eq}_2(\Sigma)=5$ in
\cite[Table~1]{AbiadJabalAmeliReijnders2025}.
\end{proof}

\begin{theorem}[Uniqueness of $q=11$]
\label{thm:why11}
Let $q$ be a prime power. If a six-arc $A\subset\PG(2,q)$ satisfies
$\cU(A)=\cC(\F_q)$ for some nonsingular conic $\cC$, then $q=11$.
At $q=11$ every such $A$ is projectively equivalent to the
Clebsch hexagon.
\end{theorem}

\begin{proof}
Theorem~\ref{cor:conic-filling-window} gives $q\ge9$ and
$c(A)=(q-6)(q-9)$ with $0\le c(A)\le15$. Thus the only prime-power
candidates are $q=9,11$: $q=10$ is not a prime power, and every
$q\ge12$ gives $c(A)\ge18$.

It remains to exclude $q=9$. Since every point of $\cC$ is uncovered, no
chord of $A$ meets $\cC$; every chord is therefore passant. A point off
$\cC$ lies on five passants if it is internal and four if it is external,
so the five chords through each vertex force all six vertices of $A$ to be
internal. Lemma~\ref{lem:q9-polarity} would then make $A$ a
six-clique in the distance-two graph of the Sylvester graph, whose clique
number is five. This is impossible.

Therefore $q=11$. The final assertion follows from
Theorem~\ref{thm:rigidity}, which puts every conic-filling six-arc over
$\F_{11}$ in the single Clebsch orbit.
\end{proof}

\section{The binary passant/internal incidence code at
\texorpdfstring{$q=13$}{q=13}}
\label{sec:q13-tangent-code}

Fix the conic $\cC:XZ-Y^2=0$ in $\PG(2,13)$.  Let $M$ be the
$78$-by-$78$ incidence matrix over $\F_2$, with rows indexed by the
passant lines to $\cC$ and columns indexed by its internal points.
Madison and Wu determine the nullity of this matrix for every odd
$q$~\cite[Corollary~6.3]{MadisonWu2012}; their formula
$(q-1)^2/4$ gives $36$ at $q=13$.  Hollmann and Xiang construct the
elliptic association scheme afforded by the
$\operatorname{PGL}(2,q)$-action on the passant carrier
\cite[\S3]{HollmannXiang2006}.  Thus the dimension and the association
scheme used below are known.  The assertions here are the exact minimum
distance and the intrinsic reconstruction of the incidence matrix and
scheme from the minimum words.

\begin{theorem}[Passant-code theorem]
\label{thm:q13-tangent-code}
The binary code $K=\ker M$ has parameters
\[
 [78,36,12]_2.
\]
There are exactly $364$ minimum words.  Under
$\operatorname{PGL}(2,13)$ they form four orbits of size $91$, one with
stabilizer $S_4$ and three with stabilizer $D_{24}$.  Each orbit spans
$K$.

The minimum supports intrinsically reconstruct the full six-class
elliptic association scheme on the $78$ internal points.  Pair concurrence
recovers passant versus secant join type, triple concurrence separates the
six orbitals, and the $78$ all-zero-triple seven-cliques are exactly the
passant incidence rows.  Consequently the minimum layer reconstructs $M$
up to row order, and its automorphism group is
$\operatorname{PGL}(2,13)$.
\end{theorem}

\begin{proof}
Every row and column of $M$ has weight seven, and conic polarity identifies
the two index sets.  Madison and Wu's dimension formula gives
$\dim K=36$, equivalently $\operatorname{rank}M=42$; exact elimination
independently checks this specialization.
Summing the row equations gives the all-one functional, so every word of
$K$ has even weight.  Moreover, if a nonzero word contains an internal
point $P$, each of the seven passants through $P$ must contain another
support point.  These seven points are distinct, proving the elementary
lower bound $\operatorname{wt}\ge8$.

Suppose first that a word has weight eight.  The preceding parity argument
shows that its support is an eight-arc, every join of two support points is
passant, and the passant pencil at every vertex is saturated.  Thus the
seven remaining lines through each vertex are precisely the seven
combinatorial tangents to the arc and the seven conic-secants through that
internal point.

For an internal point $P$, let
\[
 T_P(X)=\prod_{\substack{\ell\ni P\\
                         \ell\text{ secant to }\cC}}\ell(X).
\]
For a pairwise-passant triple put
\[
 h(P,Q,R)=
 \frac{T_P(Q)T_Q(R)T_R(P)}
      {T_P(R)T_R(Q)T_Q(P)}.
\]
Segre's lemma of tangents
\cite[Lemma~27]{BallLavrauw2019} gives $h(P,Q,R)=1$ for every triple in the
hypothetical support.  Fixing $P=(1:0:2)$ therefore makes its other seven
points a clique in the graph on the $42$ passant-join neighbors of $P$,
where $Q$ and $R$ are adjacent when $QR$ is passant and $h(P,Q,R)=1$.

An order-$14$ element of the stabilizer of $P$ splits these vertices into
three cyclic orbits $A_i,B_i,C_i$, with indices in $\mathbf Z/14$.  Direct
substitution in the displayed product gives the complete adjacency rule
\[
\begin{array}{c|c}
\text{orbit pair}&j-i\pmod {14}\\ \hline
AA&4,6,8,10\\
AB&6,7,11,12\\
AC&1,3\\
BB&6,8\\
BC&3,5,6,8,9,11\\
CC&2,4,10,12.
\end{array}
\]
Up to simultaneous translation, every four-clique is one of five rows:
\[
\begin{array}{c|c}
\text{four-clique}&\text{unique extension}\\ \hline
A_0,B_6,B_{12},C_1&C_3\\
A_0,B_6,B_{12},C_3&C_1\\
A_0,B_6,C_1,C_3&B_{12}\\
A_0,B_{12},C_1,C_3&B_6\\
B_0,B_6,C_9,C_{11}&A_8.
\end{array}
\]
They close to fourteen translates of one five-clique, none of which
extends.  The local clique number is therefore five, not seven.  This
excludes weight eight.

For weight ten, fix one support point.  Parity leaves only two passant-pencil
profiles.  If $s$ of the other nine support points are joined to the fixed
point by secants, the remaining $9-s$ points are distributed in positive
odd numbers among its seven passants.  Hence $s$ is even and
$9-s\ge7$, so $s=0$ or $2$.  These give, respectively, three points on one
passant and one on each of the other six, or one on every passant together
with two secant-join neighbors.  Exact
meet-in-the-middle syndrome checks exhaust respectively
\[
 7\binom63 6^6=6{,}531{,}840,\qquad
 6^7\binom{35}{2}=166{,}561{,}920
\]
supports, and neither profile has zero syndrome.  Finally, the following
twelve internal points do have zero syndrome:
\[
\begin{gathered}
(1:0:2),(1:3:2),(1:4:5),(1:1:8),\\
(1:4:8),(1:1:7),(1:7:12),(1:3:3),\\
(1:9:11),(1:10:11),(1:0:5),(1:8:7).
\end{gathered}
\]
Their stabilizer is dihedral of order $24$.  This proves $d(K)=12$.

The classification of the minimum layer is exhaustive and exact.  The replay
generates four projective orbits from explicit representatives and verifies
all $364$ supports against $M$.  It then forms all pair and triple concurrence
profiles.  Pair concurrences are $7,9,12$ on passant joins and $6,8$ on
secant joins; the six triple histograms are distinct.  Exactly $78$ of the
resulting passant seven-cliques have zero triple concurrence, and direct
comparison identifies them with the rows of $M$.

The span and automorphism conclusions have a shorter structural proof.  Put

\[
 \Delta(x,y,z)=y^2-xz,\qquad
 \beta(P,Q)=2y_Py_Q-x_Pz_Q-z_Px_Q,\qquad
 \rho(P,Q)=\frac{\beta(P,Q)^2}{\Delta(P)\Delta(Q)}.
\]
On distinct internal points, $\rho$ takes the six values
$0,1,3,9,10,12$.  Its level relations are the six classes of the elliptic
association scheme of Hollmann and Xiang~\cite[\S3]{HollmannXiang2006}; write
$A_r$ for the adjacency matrix of the relation $\rho=r$.  Conic polarity
identifies an internal point with its passant polar, so $A_0$ is $M$, up to
row order.

For clarity, the part of the intersection algebra used here can also be
checked directly.  Fixing one point and substituting one representative of
each $\rho$-relation gives, over $\F_2$,
\begin{equation}
\label{eq:q13-mod-two-scheme}
\begin{split}
 A_0^2&=I+A_9+A_{10}+A_{12},\\
 A_0A_r&=0\quad(r=9,10,12),\\
 A_9^2&=A_{10},\qquad A_{10}^2=A_{12},\qquad A_{12}^2=A_9.
\end{split}
\end{equation}
These are parity reductions of ordinary integer intersection numbers, not
identities inferred from matrix rank.  For example the coefficients of
$A_0^2$ on
$I,A_0,A_1,A_3,A_9,A_{10},A_{12}$ are
$(7,0,0,0,1,1,1)$; those of $A_0A_9,A_0A_{10},A_0A_{12}$ are respectively
\[
 (0,2,2,2,0,2,0),\quad
 (0,2,2,0,2,0,2),\quad
 (0,2,0,2,0,2,2),
\]
and those of $A_9^2,A_{10}^2,A_{12}^2$ are respectively
\[
 (14,0,4,2,2,1,4),\quad
 (14,0,2,4,2,4,1),\quad
 (14,4,2,2,3,2,2).
\]

Let $B=A_9$.  Since $A_0B=0$, the image of $B$ lies in
$K=\ker A_0$.  If $x\in K$ and $Bx=0$, the first identity in
\eqref{eq:q13-mod-two-scheme}, together with $B^2=A_{10}$ and
$B^4=A_{12}$, gives
\[
 0=A_0^2x=(I+B+B^2+B^4)x=x.
\]
Thus $B$ is injective on the $36$-dimensional space $K$, so
$\operatorname{im}B=K$ and $B,B^2,B^4$ all have rank $36$.  If $N_i$ is
the $91$-by-$78$ support matrix of the $i$th minimum-word orbit, direct
pair counting in one representative orbit gives
\[
 (N_1^{\mathsf T}N_1,N_2^{\mathsf T}N_2,
   N_3^{\mathsf T}N_3,N_4^{\mathsf T}N_4)
   =(A_9,A_9,A_{12},A_{10}).
\]
Hence every $N_i$ has rank at least $36$; its rows lie in $K$, so every
minimum-word orbit spans $K$.

Finally, the six concurrence colors recover exactly the six $\rho$-relations,
so an automorphism of the minimum layer preserves the elliptic scheme.  The
following four internal points form a geometric base:
\[
 P_0=(1:0:2),\quad P_1=(1:1:7),\quad
 P_2=(1:0:7),\quad P_3=(1:1:3).
\]
Their first three pair relations are
$(\rho(P_0,P_1),\rho(P_0,P_2),\rho(P_1,P_2))=(10,3,9)$.
There are $78\cdot14\cdot2=2184$ ordered triples with this relation pattern:
the relation $A_{10}$ has valency $14$, and the relevant intersection number
is $p_{3,9}^{10}=2$.  Direct substitution in the symmetric-square action
shows that the stabilizer of $(P_0,P_1,P_2)$ in
$\operatorname{PGL}(2,13)$ is trivial.  Since that group also has order
$2184$, it acts simply transitively on these triples.  For a fixed such
triple, $P_3$ is the unique internal point with relation signature
$(3,1,9)$, and the four values
\[
 \bigl(\rho(P,P_0),\rho(P,P_1),\rho(P,P_2),\rho(P,P_3)\bigr),
\]
with equality to an anchor recorded on the diagonal, distinguish all $78$
internal points.  Both assertions follow by substitution in the displayed
formula for $\rho$; the exact replay checks all $78$ normalized internal
points independently.  Therefore a scheme automorphism can first be composed
with a unique element of $\operatorname{PGL}(2,13)$ so that it fixes the first
three anchors; it then fixes the fourth anchor and hence every internal point.
Thus the reconstructed minimum layer has automorphism group exactly
$\operatorname{PGL}(2,13)$.  The original rank eliminations and the independent
order-$2184$ stabilizer-chain enumeration remain in the replay as checks.
\end{proof}

\section{The small-arc boundary}

The chord moments reduce the neighboring cases through eight points to a
finite list of fields. Exact terminal enumeration excludes the remaining
seven- and eight-point cases.

\begin{theorem}[Conic-filling loci through eight points]
\label{thm:small-k-conic-filling}
Let $A$ be a $k$-arc in $\PG(2,q)$, where $q$ is a prime power and
$4\le k\le8$, and suppose that $\cU(A)$ is the full set of
$\F_q$-rational points of a nonsingular conic. Then
\[
 (k,q)=(4,5)\qquad\text{or}\qquad(k,q)=(6,11).
\]
In the first case $A$ is a projective four-frame; in the second it is a
Clebsch hexagon. Conversely, both cases occur.

Moreover, for each $q\in\{13,17,19\}$, an arc all of whose chords are
passant to a fixed nonsingular conic has at most six points, and this bound
is attained.
\end{theorem}

\begin{proof}[Proof with exhaustive finite verification]
For $k=4$, Theorem~\ref{cor:conic-filling-window} gives
$5\le q<6$, hence $q=5$. Every four-arc is a projective frame, and for the
standard frame a direct substitution gives
\[
 \cU(A)=V(X^2+Y^2+Z^2+XY+XZ+YZ),
\]
where $V(f)$ denotes the projective zero set of $f$; this is a nonsingular
six-point conic. For $k=5$, Theorem~\ref{lem:chord-defect} gives
$|\cU(A)|=q^2-9q+21$; equality with $q+1$ would require
$q^2-10q+20=0$, which has no integral root. The case $k=6$ is
Theorem~\ref{thm:why11} together with Theorem~\ref{thm:rigidity}.

For $k=7$, Theorem~\ref{lem:chord-defect} gives
\[
 |\cU(A)|=q^2-20q+120-n_3.
\]
If this equals $q+1$, then
\[
 n_3=q^2-21q+119,\quad
n_2=105-3n_3,\quad
 n_1=3(q^2-14q+42).
\]
The strengthened conic-filling window gives $11\le q\le15$. Thus only the
prime powers $q=11,13$ remain, with forced spectra
$(n_1,n_2,n_3)=(27,78,9)$ and $(87,60,15)$ respectively. An exhaustive
four-frame-normalized enumeration now supplies the finite exclusion. It
finds $10232$ distinct normalized seven-arcs at $q=11$, of which $140$ have
$|\cU(A)|=q+1$, and $53960$ at $q=13$, of which $1680$ have that size.
Every seven-arc is reached because it contains a four-frame; extending each
normalized six-arc by its uncovered points produces every resulting
seven-arc exactly three times, according to which non-frame point is
deleted. For all $1820$ candidates of the required size, the quadratic
evaluation matrix has full rank six, so none lies on a conic.

Finally let $k=8$. The strengthened conic-filling window gives
$13\le q\le59/3$, so the odd prime powers are exactly $13,17,19$.

It remains to exclude these three fields. By projective equivalence fix the
conic $XZ=Y^2$, whose stabilizer is the symmetric-square image of
$\operatorname{PGL}(2,q)$. Its complement has $q^2$ points. In the
terminology of Blokhuis, Seress, and Wilbrink, a set whose pairwise joins
are passant is an exterior set of the conic
\cite[pp.~143,146]{BSW1992}; its vertices may have either internal or
exterior point type. Every conic-filling arc is such an exterior set.

There is already a structural reduction at $q=13$. An exterior point lies
on only $(q-1)/2=6$ passants, fewer than the seven distinct chords through
a vertex of an eight-arc. Hence every vertex would be internal. Each
internal point lies on exactly $(q+1)/2=7$ passants, so its seven chords
would exhaust its complete passant pencil.  Its characteristic vector
would therefore be a weight-eight word of the code in
Theorem~\ref{thm:q13-tangent-code}, which is impossible.  This gives the
conceptual $q=13$ exclusion needed for conic filling.

The stronger assertion that every pairwise-passant arc has at most six
points still needs the orbit search.  Join two off-conic points when their
line is passant. The desired eight-arc would be an eight-clique in this
graph, with collinear triples forbidden.

An exact orbit search gives the following data:
\[
\begin{array}{c|c|c|c}
q&\text{off-conic points}&\text{passant edges}&
 \operatorname{PGL}(2,q)\text{-orbits on edges}\\ \hline
13&169&7098&10\\
17&289&20808&13\\
19&361&32490&15.
\end{array}
\]
For one representative of each edge orbit, the search extends partial arcs
and identifies states under the setwise stabilizer of the root edge.
Canonical keys are the lexicographically least sorted point lists in each
stabilizer orbit. Every candidate is reached: it contains a passant edge,
that edge is carried to one of the listed representatives, and the remaining
points are then among the enumerated extensions. In all three fields the
largest partial arc has size six. Thus there is no seven-point arc with all
chords passant, and in particular no conic-filling eight-arc. A separate
replay verifies the edge-orbit partition and repeats the extension search by
ordered backtracking without stabilizer identification. Its certificates
give an explicit six-point witness in every field, proving sharpness.
\end{proof}

In coding terms, let $C$ be a projective $[k,k-3,4]_q$ MDS code with
$4\le k\le8$. If its projective distance-three syndrome locus is a
nonsingular conic, then, up to monomial equivalence, $C$ is the
$[4,1,4]_5$ frame code or the $[6,3,4]_{11}$ Clebsch code. Both codes have
the stated syndrome locus.

Thus, for each fixed arc size, only finitely many field orders can admit
conic filling; the classification is complete through eight points.

\section{Trust and exact replay}
\label{sec:verification}

The claims in this paper use three distinct proof modes.  The chord-defect,
window, and rigidity inputs are proved in the geometric paper.  The
$q=9$ graph identification and clique bound use the cited literature, as
does the tangent identity in the conceptual $q=13$ reduction.
The $q=13$ orbit-span and automorphism conclusions are proved structurally
here from the elliptic intersection algebra and a four-point geometric base;
the replay checks those arguments independently.  The census, minimum-layer
classification, weight-profile exclusions, and terminal orbit searches are
exhaustive finite computations.  Table~\ref{tab:replays} lists the load-bearing
executable route for each such claim; these checks use exact finite-field
arithmetic and integer combinatorics, not randomized tests.

The machine-readable version of this ledger is the file
\begin{center}
\small\path{verification/computational_companion_trust.json}.
\end{center}
Its paths are relative to the present directory.  Recorded outputs and
source hashes are in \path{verification/checker_outputs.json}; the larger
$k=8$ certificates are checked by
\path{verification/conic_filling_verify.py}.  The clean companion entry point is
\begin{center}
\small\path{verification/verify_computational_companion.py}.
\end{center}
Passing \path{--run} also executes every replay in
Table~\ref{tab:replays}.

\begin{table}[H]
\centering
\footnotesize
\begin{tabular}{@{}>{\raggedright\arraybackslash}p{0.27\textwidth}
                    >{\raggedright\arraybackslash}p{0.65\textwidth}@{}}
\toprule
Result & Exhaustive or independent executable check \\
\midrule
Fifteen classes and numerical gap
& \path{check_global_conic_gap.py} \\
Low-degree rigidity and quartic witness
& \path{check_low_degree_loci.py} \\
$q=9$ exclusion and uniqueness of $q=11$
& \path{check_small_q_uniqueness.py} \\
Seven-arc exclusions at $q=11,13$
& \path{check_small_k_conic_filling.py} \\
$q=13$ passant/internal incidence code and minimum layer
& \path{check_q13_tangent_code.py} \\
Passant-arc bound for $q=13,17,19$
& \path{verification/conic_filling_verify.py} \\
\bottomrule
\end{tabular}
\caption{Load-bearing exact replays for the computational claims.}
\label{tab:replays}
\end{table}

No exhaustive classification in this paper is certified by Lean. A
separate development checks the chord moments, six-arc defect
specialization, the $q=9$ polarity reduction, and the geometric bridge used
in the small-$k$ argument. Accordingly, the replays in
Table~\ref{tab:replays} are part of the proof surface rather than optional
corroboration.

\section{Conclusion}

The finite census exposes two strengthenings that the geometric proof does
not need: degree-three containment already characterizes the Clebsch class,
and uncovered-set size has a four-point gap above it.  Across fields, the
Sylvester obstruction makes $q=11$ unique for conic-filling six-arcs.
Together with the terminal searches, this leaves exactly two
conic-filling projective codes through length eight.  At the first
terminal field, the passant-code argument also identifies the exact binary
distance and shows that the complete minimum layer reconstructs its
passant geometry.

\section*{AI assistance disclosure}

This project was developed with extensive assistance from OpenAI Codex and
Anthropic Claude. Under the author's direction, the systems assisted
throughout with proof exploration, literature research, symbolic and finite
computation, code and formal-proof development, verification, and manuscript
drafting and revision. The author checked the resulting arguments,
computations, code, and cited sources, reviewed and edited all AI-assisted
material, and assumes responsibility for all content.

\begin{thebibliography}{9}

\bibitem{AbiadJabalAmeliReijnders2025}
A.~Abiad, A.~Jabal Ameli, and L.~Reijnders,
\newblock The clique number of the exact distance $t$-power graph:
complexity and eigenvalue bounds,
\newblock \emph{Discrete Applied Mathematics} \textbf{363} (2025), 55--70.
\newblock doi:10.1016/j.dam.2024.11.032.

\bibitem{BallLavrauw2019}
S.~Ball and M.~Lavrauw,
\newblock Arcs in finite projective spaces,
\newblock \emph{EMS Surveys in Mathematical Sciences} \textbf{6} (2019),
133--172.
\newblock arXiv:1908.10772.

\bibitem{BSW1992}
A.~Blokhuis, Á.~Seress, and H.~A.~Wilbrink,
\newblock Characterization of complete exterior sets of conics,
\newblock \emph{Combinatorica} \textbf{12}(2) (1992), 143--147.
\newblock doi:10.1007/BF01204717.

\bibitem{BrouwerCohenNeumaier1989}
A.~E.~Brouwer, A.~M.~Cohen, and A.~Neumaier,
\newblock \emph{Distance-Regular Graphs},
\newblock Ergebnisse der Mathematik und ihrer Grenzgebiete, vol.~18,
Springer-Verlag, Berlin, 1989.

\bibitem{Dye1991}
R.~H.~Dye,
\newblock Hexagons, conics, $A_5$ and $\mathrm{PSL}_2(K)$,
\newblock \emph{Journal of the London Mathematical Society} (2) \textbf{44}
(1991), 270--286.
\newblock doi:10.1112/jlms/s2-44.2.270.

\bibitem{HollmannXiang2006}
H.~D.~L.~Hollmann and Q.~Xiang,
\newblock Association schemes from the action of
$\mathrm{PGL}(2,q)$ fixing a nonsingular conic in $\PG(2,q)$,
\newblock \emph{Journal of Algebraic Combinatorics} \textbf{24}(2)
(2006), 157--193.
\newblock doi:10.1007/s10801-006-0005-8; arXiv:math/0503573.

\bibitem{JurisicVidali2019}
A.~Juri\v{s}i\'c and J.~Vidali,
\newblock The Sylvester graph and Moore graphs,
\newblock \emph{European Journal of Combinatorics} \textbf{80} (2019),
184--193.
\newblock doi:10.1016/j.ejc.2018.02.018.

\bibitem{MadisonWu2012}
A.~L.~Madison and J.~Wu,
\newblock On binary codes from conics in $\PG(2,q)$,
\newblock \emph{European Journal of Combinatorics} \textbf{33}(1)
(2012), 33--48.
\newblock doi:10.1016/j.ejc.2011.08.001; arXiv:1104.0324.

\bibitem{RuddRigidity2026}
T.~Rudd,
\newblock Reconstructing the Clebsch code and its golden orientation from
its deep-hole syndrome locus,
\newblock companion manuscript, 2026.

\end{thebibliography}

\end{document}
